%
% Standard operations from linear algebra.
%
\ifx\havemjolinearalgebra\undefined
\def\havemjolinearalgebra{1}


\ifx\lvert\undefined
  \usepackage{amsmath} % \lvert, \rVert, etc. and \operatorname.
\fi

\ifx\ocircle\undefined
  \usepackage{wasysym}
\fi

\ifx\clipbox\undefined
  % Part of the adjustbox package; needed to clip the \perp sign.
  \usepackage{trimclip}
\fi

%
% Only the most commonly-used macros. Needed by everything else.
%
\ifx\havemjocommon\undefined
\def\havemjocommon{1}

\ifx\mathbb\undefined
  \usepackage{amsfonts}
\fi

\ifx\restriction\undefined
  \usepackage{amssymb}
\fi

% Place the argument in matching left/right parentheses.
\newcommand*{\of}[1]{ \left({#1}\right) }

% Group terms using parentheses.
\newcommand*{\qty}[1]{ \left({#1}\right) }

% Group terms using square brackets.
\newcommand*{\sqty}[1]{ \left[{#1}\right] }

% A pair of things.
\newcommand*{\pair}[2]{ \left({#1},{#2}\right) }

% A triple of things.
\newcommand*{\triple}[3]{ \left({#1},{#2},{#3}\right) }

% A four-tuple of things.
\newcommand*{\quadruple}[4]{ \left({#1},{#2},{#3},{#4}\right) }

% A five-tuple of things.
\newcommand*{\quintuple}[5]{ \left({#1},{#2},{#3},{#4},{#5}\right) }

% A six-tuple of things.
\newcommand*{\sextuple}[6]{ \left({#1},{#2},{#3},{#4},{#5},{#6}\right) }

% A seven-tuple of things.
\newcommand*{\septuple}[7]{ \left({#1},{#2},{#3},{#4},{#5},{#6},{#7}\right) }

% A free-form tuple of things. Useful for when the exact number is not
% known, such as when \ldots will be stuck in the middle of the list,
% and when you don't want to think in column-vector terms, e.g. with
% elements of an abstract Cartesian product space.
\newcommand*{\tuple}[1]{ \left({#1}\right) }

% The "least common multiple of" function. Takes a nonempty set of
% things that can be multiplied and ordered as its argument. Name
% chosen for synergy with \gcd, which *does* exist already.
\newcommand*{\lcm}[1]{ \operatorname{lcm}\of{{#1}} }
\ifdefined\newglossaryentry
  \newglossaryentry{lcm}{
    name={\ensuremath{\lcm{X}}},
    description={the least common multiple of the elements of $X$},
    sort=l
  }
\fi

% The factorial operator.
\newcommand*{\factorial}[1]{ {#1}! }

% Restrict the first argument (a function) to the second argument (a
% subset of that functions domain). Abused for polynomials to specify
% an associated function with a particular domain (also its codomain,
% in the case of univariate polynomials).
\newcommand*{\restrict}[2]{{#1}{\restriction}_{#2}}
\ifdefined\newglossaryentry
  \newglossaryentry{restriction}{
    name={\ensuremath{\restrict{f}{X}}},
    description={the restriction of $f$ to $X$},
    sort=r
  }
\fi

%
% Product spaces
%
% All of the product spaces (for example, R^n) that follow default to
% an exponent of ``n'', but that exponent can be changed by providing
% it as an optional argument. If the exponent given is ``1'', then it
% will be omitted entirely.
%

% The natural n-space, N x N x N x ... x N.
\newcommand*{\Nn}[1][n]{
  \mathbb{N}\if\detokenize{#1}\detokenize{1}{}\else^{#1}\fi
}

\ifdefined\newglossaryentry
  \newglossaryentry{N}{
    name={\ensuremath{\Nn[1]}},
    description={the set of natural numbers},
    sort=N
  }
\fi

% The integral n-space, Z x Z x Z x ... x Z.
\newcommand*{\Zn}[1][n]{
  \mathbb{Z}\if\detokenize{#1}\detokenize{1}{}\else^{#1}\fi
}

\ifdefined\newglossaryentry
  \newglossaryentry{Z}{
    name={\ensuremath{\Zn[1]}},
    description={the ring of integers},
    sort=Z
  }
\fi

% The rational n-space, Q x Q x Q x ... x Q.
\newcommand*{\Qn}[1][n]{
  \mathbb{Q}\if\detokenize{#1}\detokenize{1}{}\else^{#1}\fi
}

\ifdefined\newglossaryentry
  \newglossaryentry{Q}{
    name={\ensuremath{\Qn[1]}},
    description={the field of rational numbers},
    sort=Q
  }
\fi

% The real n-space, R x R x R x ... x R.
\newcommand*{\Rn}[1][n]{
  \mathbb{R}\if\detokenize{#1}\detokenize{1}{}\else^{#1}\fi
}

\ifdefined\newglossaryentry
  \newglossaryentry{R}{
    name={\ensuremath{\Rn[1]}},
    description={the field of real numbers},
    sort=R
  }
\fi


% The complex n-space, C x C x C x ... x C.
\newcommand*{\Cn}[1][n]{
  \mathbb{C}\if\detokenize{#1}\detokenize{1}{}\else^{#1}\fi
}

\ifdefined\newglossaryentry
  \newglossaryentry{C}{
    name={\ensuremath{\Cn[1]}},
    description={the field of complex numbers},
    sort=C
  }
\fi

% The n-dimensional product space of a generic field F.
\newcommand*{\Fn}[1][n]{
  \mathbb{F}\if\detokenize{#1}\detokenize{1}{}\else^{#1}\fi
}

\ifdefined\newglossaryentry
  \newglossaryentry{F}{
    name={\ensuremath{\Fn[1]}},
    description={a generic field},
    sort=F
  }
\fi


% An indexed arbitrary binary operation such as the union or
% intersection of an infinite number of sets. The first argument is
% the operator symbol to use, such as \cup for a union. The second
% argument is the lower index, for example k=1. The third argument is
% the upper index, such as \infty. Finally the fourth argument should
% contain the things (e.g. indexed sets) to be operated on.
\newcommand*{\binopmany}[4]{
  \mathchoice{ \underset{#2}{\overset{#3}{#1}}{#4} }
             { {#1}_{#2}^{#3}{#4} }
             { {#1}_{#2}^{#3}{#4} }
             { {#1}_{#2}^{#3}{#4} }
}


% The four standard (UNLESS YOU'RE FRENCH) types of intervals along
% the real line.
\newcommand*{\intervaloo}[2]{ \left({#1},{#2}\right) } % open-open
\newcommand*{\intervaloc}[2]{ \left({#1},{#2}\right] } % open-closed
\newcommand*{\intervalco}[2]{ \left[{#1},{#2}\right) } % closed-open
\newcommand*{\intervalcc}[2]{ \left[{#1},{#2}\right] } % closed-closed


\fi
 % for \of, at least

% Absolute value (modulus) of a scalar.
\newcommand*{\abs}[1]{\left\lvert{#1}\right\rvert}

% Norm of a vector.
\newcommand*{\norm}[1]{\left\lVert{#1}\right\rVert}

% The inner product between its two arguments.
\newcommand*{\ip}[2]{\left\langle{#1},{#2}\right\rangle}

% The tensor product of its two arguments.
\newcommand*{\tp}[2]{ {#1}\otimes{#2} }

% The Kronecker product of its two arguments. The usual notation for
% this is the same as the tensor product notation used for \tp, but
% that leads to confusion because the two definitions may not agree.
\newcommand*{\kp}[2]{ {#1}\odot{#2} }

% The adjoint of a linear operator.
\newcommand*{\adjoint}[1]{ #1^{*} }

% The ``transpose'' of a linear operator; namely, the adjoint, but
% specialized to real matrices.
\newcommand*{\transpose}[1]{ #1^{T} }

% The Moore-Penrose (or any other, I guess) pseudo-inverse of its
% sole argument.
\newcommand*{\pseudoinverse}[1]{ #1^{+} }

% The trace of an operator.
\newcommand*{\trace}[1]{ \operatorname{trace}\of{{#1}} }

% The diagonal matrix whose only nonzero entries are on the diagonal
% and are given by our argument. The argument should therefore be a
% vector or tuple of entries, by convention going from the top-left to
% the bottom-right of the matrix.
\newcommand*{\diag}[1]{\operatorname{diag}\of{{#1}}}

% The "rank" of its argument, which is context-dependent. It can mean
% any or all of,
%
%   * the rank of a matrix,
%   * the rank of a power-associative algebra (particularly an EJA),
%   * the rank of an element in a Euclidean Jordan algebra.
%
\newcommand*{\rank}[1]{ \operatorname{rank}\of{{#1}} }


% The ``span of'' operator. The name \span is already taken.
\newcommand*{\spanof}[1]{ \operatorname{span}\of{{#1}} }

% The ``co-dimension of'' operator.
\newcommand*{\codim}{ \operatorname{codim} }

% The orthogonal projection of its second argument onto the first.
\newcommand*{\proj}[2] { \operatorname{proj}\of{#1, #2} }

% The set of all eigenvalues of its argument, which should be either a
% matrix or a linear operator. The sigma notation was chosen instead
% of lambda so that lambda can be reserved to denote the ordered tuple
% (largest to smallest) of eigenvalues.
\newcommand*{\spectrum}[1]{\sigma\of{{#1}}}
\ifdefined\newglossaryentry
  \newglossaryentry{spectrum}{
    name={\ensuremath{\spectrum{L}}},
    description={the set of all eigenvalues of $L$},
    sort=s
  }
\fi

% The reduced row-echelon form of its argument, a matrix.
\newcommand*{\rref}[1]{\operatorname{rref}\of{#1}}
\ifdefined\newglossaryentry
  \newglossaryentry{rref}{
    name={\ensuremath{\rref{A}}},
    description={the reduced row-echelon form of $A$},
    sort=r
  }
\fi

% The ``Automorphism group of'' operator.
\newcommand*{\Aut}[1]{ \operatorname{Aut}\of{{#1}} }

% The ``Lie algebra of'' operator.
\newcommand*{\Lie}[1]{ \operatorname{Lie}\of{{#1}} }

% The ``write a matrix as a big vector'' operator.
\newcommand*{\vectorize}[1]{ \operatorname{vec}\of{{#1}} }

% The ``write a big vector as a matrix'' operator.
\newcommand*{\matricize}[1]{ \operatorname{mat}\of{{#1}} }

% An inline column vector, with parentheses and a transpose operator.
\newcommand*{\colvec}[1]{ \transpose{\left({#1}\right)} }

% Bounded linear operators on some space. The required argument is the
% domain of those operators, and the optional argument is the
% codomain. If the optional argument is omitted, the required argument
% is used for both.
\newcommand*{\boundedops}[2][]{
  \mathcal{B}\of{ {#2}
    \if\relax\detokenize{#1}\relax
      {}%
    \else
      {,{#1}}%
    \fi
  }
}


%
% Orthogonal direct sum.
%
% First declare my ``perp in a circle'' operator, which is meant to be
% like an \obot or an \operp except has the correct weight circle. It's
% achieved by overlaying an \ocircle with a \perp, but only after we
% clip off the top half of the \perp sign and shift it up.
\DeclareMathOperator{\oplusperp}{\mathbin{
  \ooalign{
    $\ocircle$\cr
    \raisebox{0.625\height}{$\clipbox{0pt 0pt 0pt 0.5\height}{$\perp$}$}\cr
  }
}}

% Now declare an orthogonal direct sum in terms of \oplusperp.
\newcommand*{\directsumperp}[2]{ {#1}\oplusperp{#2} }


% The space of real symmetric n-by-n matrices. Does not reduce to
% merely "S" when n=1 since S^{n} does not mean an n-fold cartesian
% product of S^{1}.
\newcommand*{\Sn}[1][n]{ \mathcal{S}^{#1} }
\ifdefined\newglossaryentry
  \newglossaryentry{Sn}{
    name={\ensuremath{\Sn}},
    description={the set of $n$-by-$n$ real symmetric matrices},
    sort=Sn
  }
\fi

% The space of complex Hermitian n-by-n matrices. Does not reduce to
% merely "H" when n=1 since H^{n} does not mean an n-fold cartesian
% product of H^{1}. The field may also be given rather than assumed
% to be complex; for example \Hn[3]\of{\mathbb{O}} might denote the
% 3-by-3 Hermitian matrices with octonion entries.
\newcommand*{\Hn}[1][n]{ \mathcal{H}^{#1} }
\ifdefined\newglossaryentry
  \newglossaryentry{Hn}{
    name={\ensuremath{\Hn}},
    description={the set of $n$-by-$n$ complex Hermitian matrices},
    sort=Hn
  }
\fi


% The general linear group of square matrices whose size is the first
% argument and whose entries come from the second argument.
\newcommand*{\GL}[2]{\operatorname{GL}_{#1}\of{#2}}
\ifdefined\newglossaryentry
  \newglossaryentry{GL}{
    name={\ensuremath{\GL{n}{\mathbb{F}}}},
    description={the general linear group of order $n$ over $\mathbb{F}$},
    sort=GL
  }
\fi


% The space of all (linear) derivations on its argument, an algebra.
\newcommand*{\Der}[1]{\operatorname{Der}\of{#1}}
\ifdefined\newglossaryentry
  \newglossaryentry{Der}{
    name={\ensuremath{\Der{V}}},
    description={the space of all (linear) derivations on $V$},
    sort=Der
  }
\fi


% The set of all isometries from its first argument to its second
% (optional) argument, both assumed to be at least normed spaces, and
% more likely Hilbert spaces. The norms are implicit, i.e. not
% included in the arguments. If the optional argument is omitted, you
% get the isometries from the first argument to itself AKA its
% isometry group.
\newcommand*{\Isom}[2][]{
  \operatorname{Isom}\of{ {#2}
    \if\relax\detokenize{#1}\relax
      {}%
    \else
      {,{#1}}%
    \fi
  }
}
\ifdefined\newglossaryentry
  \newglossaryentry{Isom}{
    name={\ensuremath{\Isom{V}}},
    description={the group of isometries on $V$},
    sort=Isom
  }
  \newglossaryentry{Isom2}{
    name={\ensuremath{\Isom[W]{V}}},
    description={the set of isometries from $V$ to $W$},
    sort=Isom
  }
\fi

\fi
